\section{Introduction}
\label{sec:introduction}

Semidefinite lifts describe a convex set as the image of an affine
section of a positive semidefinite cone.  They can be substantially
smaller than linear descriptions, so lower bounds for linear extended
formulations do not automatically exclude small semidefinite lifts.
We prove an exponential lower bound for perfect matching by studying
an exact factorization problem that also allows a fixed positive
shift of every odd-cut slack.

For even $n\ge4$, let $E_n$ be the edges of the complete graph $K_n$
on $[n]=\{1,\ldots,n\}$, let $\mathcal M_n$ be its perfect matchings,
and let $\mathcal O_n$ be the odd-cardinality subsets of $[n]$.
For $U\subseteq[n]$, write $\delta(U)$ for the edges with exactly
one endpoint in $U$.  For $0\le\rho\le1$, define the matrix
\begin{equation}\label{eq:main-matrix}
 A_n(\rho)_{U,M}=|M\cap\delta(U)|-1+\rho,
 \qquad U\in\mathcal O_n,\quad M\in\mathcal M_n.
\end{equation}
Every perfect matching crosses an odd set an odd positive number
of times, so this matrix is nonnegative.  Its rows consist solely
of odd cuts.

The \emph{real positive semidefinite rank} of a finite nonnegative
matrix $A$, denoted $\rankpsd A$, is the least integer $r\ge1$ for
which there are real symmetric positive semidefinite matrices
$F_U,G_M\in\R^{r\times r}$ satisfying
\[
 A_{U,M}=\tr(F_UG_M)
 \qquad\text{for every row and column}.
\]
Such matrices are called PSD factors.  The entries must be reproduced
exactly, and no symmetry, equivariance, sparsity, or precision
restriction is imposed on the factors.

\begin{theorem}\label{thm:main}
There is a constant $c>0$ such that, for every fixed $0<\rho<1$,
there is $n_0=n_0(\rho)$ for which
\[
 \rankpsd A_n(\rho)\ge2^{cn}
\]
for every even $n\ge n_0$.
\end{theorem}

The endpoint $\rho=1$ has a factorization of size $\binom n2$:
use diagonal matrices indexed by edges, with diagonal entries
$\mathbf1_{\{e\in\delta(U)\}}$ for a row and
$\mathbf1_{\{e\in M\}}$ for a column.  Thus the endpoint $\rho=1$
cannot be included in the exponential lower bound.
The exponential rate in Theorem~\ref{thm:main} is independent of
$\rho$, although the threshold is allowed to deteriorate as
$\rho$ tends to one.

\subsection{Exact semidefinite lifts}

Let
\[
 P_n=\operatorname{conv}\{\mathbf1_M:M\in\mathcal M_n\}
 \subseteq\R^{E_n}
\]
be the perfect matching polytope.  A \emph{semidefinite lift of
size $r$} represents $P_n$ as an affine image of
$L\cap\mathbb S_+^r$, where $L$ is an affine subspace of the real
symmetric matrices and $\mathbb S_+^r$ is their PSD cone.  Also let
$T_n(\rho)$ denote the matrix obtained by adjoining to $A_n(\rho)$
the edge rows
\[
 T_n(\rho)_{e,M}=\mathbf1_{\{e\in M\}},\qquad e\in E_n.
\]
Adding rows cannot decrease PSD rank, so
Theorem~\ref{thm:main} applies to $T_n(\rho)$ as well.

\begin{corollary}[Unshifted rank and exact lifts]
\label{cor:exact-lift}
There are constants $c_0>0$ and $n_1$ such that, for every even
$n\ge n_1$, both $A_n(0)$ and $T_n(0)$ have real PSD rank at least
$2^{c_0n}$, and every exact semidefinite lift of $P_n$ has size at
least $2^{c_0n}$.
\end{corollary}
\begin{proof}
Appending scalar blocks $\rho$ to the row factors and $1$ to the
column factors of $A_n(0)$ gives a factorization of $A_n(\rho)$.
Consequently
\[
 \rankpsd A_n(\rho)\le\rankpsd A_n(0)+1.
\]
Fix $\rho=1/2$ and apply Theorem~\ref{thm:main}.  For all sufficiently
large even $n$,
\[
 \rankpsd T_n(0)\ge\rankpsd A_n(0)\ge2^{cn}-1\ge2^{cn/2}.
\]

Edmonds's description of $P_n$ uses nonnegativity and odd-cut
inequalities, within the affine space of degree equations
$\sum_{e\ni v}x_e=1$ \cite{edmonds}.  Thus $T_n(0)$ is a slack
matrix of this description.  Every row inequality is tight at some
perfect matching.  Any redundant row is therefore a nonnegative
combination of facet slack rows, so including it does not change
PSD rank.  The cone-factorization theorem identifies this rank with
the minimum size of a semidefinite lift
\cite[Theorem~2.4 and Corollary~2.6]{gpt}.

For completeness, an arbitrary lift can be restricted to the
smallest face of the PSD cone containing its feasible section.
This face is a PSD cone of size at most $r$, and the section meets
its relative interior, supplying the properness needed for
factorization.  An affine image can be made linear on the affine
slice without increasing size: the slice cannot contain zero,
since its intersection with the cone would then have a bounded
affine image only if that image were a point.  A linear functional
equal to one on the slice therefore absorbs the translation.
Taking $c_0=c/2$ proves the corollary.
\end{proof}

\subsection{History and the role of the shift}

Edmonds's inequality description \cite{edmonds} is the starting
point for formulation complexity of matching.  Yannakakis
developed the nonnegative-factorization framework and proved an
exponential lower bound for symmetric linear formulations
\cite{yannakakis}.  Rothvo\ss\ proved the unrestricted exponential
linear extension-complexity bound \cite{rothvoss}.  The
cone-factorization framework of Gouveia, Parrilo, and Thomas
\cite{gpt} gives the corresponding connection between PSD rank
and semidefinite lifts.

For the particular shifted matrix, Braun and Pokutta proved
exponential nonnegative rank for the entries
$(|M\cap\delta(U)|-1+\rho)/2$, for each fixed $0<\rho<1$
\cite[Section~2.2.3 and Theorem~3.1]{braun-pokutta}.
Positive scalar multiplication does not change rank, so their
matrix is the same factorization problem with nonnegative scalar
factors in place of arbitrary PSD factors.  Braun and collaborators
proved exponential lower bounds for coordinate-symmetric
semidefinite formulations of matching
\cite[Definition~2.2 and Theorem~3.10]{matching-symmetric}.
Here no group action on the factor coordinates is assumed.
Permutation symmetry will be used only to analyze averaging
operators.

Lee, Raghavendra, and Steurer developed a general approach to
unrestricted SDP lower bounds through sums of squares and
low-degree separating functionals \cite{lrs}.  Our proof follows
that approach, with a transfer argument tailored to matching.
It treats the exact positive shift in Equation~\eqref{eq:main-matrix}.
Exactness matters: Kaniewski, Lee, and de Wolf construct matrices
of subexponential PSD rank that approximate the unshifted matching
slack with entrywise error at most $2^{-\sqrt{n/2}}$
\cite[Section~7.3 and Theorem~19]{kaniewski-lee-dewolf}.

\subsection{Proof strategy and reusable ingredients}

The obstruction is an inconsistent parity system on an expanding
graph $H$ with $t$ vertices.  Each vertex carries $d$ incidence
bits of prescribed parity, and the sum of the prescribed parities
is odd.  Some edge must therefore have unequal bits at its two
ends.  Subtracting $1-\rho$ from the number of disagreements gives
a positive function $f$.  Products of signs of the free incidence
bits give Fourier characters, each involving certain vertex blocks.
Section~\ref{sec:parity} constructs a linear functional that is
nonnegative on squares of functions whose Fourier characters each
involve at most a small fixed fraction of the blocks, yet takes
value $-(1-\rho)$ on $f$.  The signed
moment construction is related to Grigoriev's low-degree
functionals for Tseitin contradictions \cite{grigoriev}.

To realize this system in matching, replace every vertex of $H$
by a cell of $k$ vertices.  The cell size $k$ is a sufficiently
large constant, fixed before $n$ tends to infinity.  Each auxiliary
incidence uses one terminal; the remaining vertices are paired
inside the cell.  A local cut restriction fixes the terminal bits
and gives equal bits to every internal pair.  The only crossing
edges are then the auxiliary disagreements.  Section~\ref{sec:matching}
performs this construction and completes the proof.

The main analytic task is to show that a small PSD factorization
would make $f$ close enough to a low-degree sum of squares.
Let $\overline F(y)$ be the row factor averaged over cuts with
incidence assignment $y$.  Pairing this average with the factor of
the corresponding matching gives $f(y)$.  Taking square roots of
the row factors gives matrix functions with no initial degree bound;
averaging their squares cannot be replaced by squaring their
averages.  We instead construct an exact representation
\[
 \overline F(y)=\sum_q H_q(y)^{\mathsf T}H_q(y)
\]
by real rectangular matrix functions.  The index $q$ labels both
a summand, called a branch, and a Fourier character $\chi_q$.
The goal is to concentrate the coefficients of $H_q$ on characters
that differ from $\chi_q$ in only a few vertex blocks.

Section~\ref{sec:local} proves the local averaging estimate used
at each step.  Its scalar variance bound tends to zero as $k$
grows and extends to matrices without a dimension factor.
The proof combines a trace estimate with a lower bound on the
dimensions of invariant eigenspaces, a spectral mechanism also
used in quasirandom-group arguments \cite{gowers}.
Section~\ref{sec:splitting} proves a dimension-free orthogonal
splitting inequality for a finite family of PSD matrices.

These ingredients yield the abstract product averaging theorem
in Section~\ref{sec:fourier}.  It bounds the Fourier mass of each
branch with weights growing exponentially in the number of blocks
where a character differs from its center $q$.  The controlled quantity is
the squared Frobenius norm of a weighted Gram matrix, hence a
fourth-order expression in the Fourier coefficients.  Its use
allows repeated averaging even though representing an average by
a direct sum introduces additional rows into the matrices $H_q$.
After truncation near the center $q$, multiplication by $\chi_q$
moves each retained character to low block degree while
preserving squares, as in finite-group Fourier sums of squares
\cite{fsp-sparse}.  This is the step that lets the separating
functional contradict the factorization.

The local averaging lemma, the orthogonal splitting inequality,
and the product theorem are stated separately from the global
matching factorization problem.
Appendix~\ref{sec:matching-smoothing} gives a further independent
smoothing theorem for averages over matchings with any prescribed
edge-count table across a fixed number of cells.  Its estimate is
uniform even for rare tables and unbalanced cell sizes, and its
three-cell specialization includes a single polynomial in the
three intercell counts.

\paragraph{Notation.}
Finite sets carry uniform probability unless another law is specified;
averages are denoted by $\E$.  For real matrices we use the Frobenius
norm $\normF B=(\tr(B^{\mathsf T}B))^{1/2}$ and operator norm
$\normop B$.  The relation $A\preceq B$ means that $B-A$ is PSD.
All logarithms are natural unless their base is displayed.
